(1) 在非简并、热和化学平衡下，化学势满足 $\mu_{H^0}=\mu_p+\mu_e$。用理想气体平动配分函数与电离能的Boltzmann因子，可写Saha关系 $n_pn_e/n_{H^0}=S(T)$，其中题设的 $S(T)$ 是已知正函数。
(2) 核数守恒与电中性给 $n_{H^0}+n_p=n_H$、$n_e=n_p$。若定义 $x=n_p/n_H$，则 $n_p=n_e=n_Hx$、$n_{H^0}=n_H(1-x)$。代入平衡式得 $n_Hx^2/(1-x)=S$，亦即 $x^2+sx-s=0$，其中 $s=S/n_H$。
(3) 物理解在 $0<x<1$，为 $x=(\sqrt{s^2+4s}-s)/2$，也可写成稳定计算式 $x=2s/(\sqrt{s^2+4s}+s)$。当 $s\ll1$，$x\simeq\sqrt s$；当 $s\gg1$，$1-x\simeq1/s$，符合弱电离与近完全电离两端的趋势。
(4) 电子分压为 $p_e=n_Hxk_BT$；将中性氢、离子和电子全部计数，总压强为 $p=n_H(1+x)k_BT$。固定 $T$ 降低总氢核密度使 $s=S/n_H$ 增大，电离度 $x$ 因而增大；高密度非理想和电子简并时不能直接使用此理想式。
